\chapter{Conclusão}
\label{conclusao}

\epigraph{\textit{"Vou viver as coisas novas que também são boas,\\
o amor, o humor das praças, cheias de pessoas.\\
Agora eu quero tudo, tudo outra vez..."}}{\textbf{--- Belchior}}

Este trabalho teve como propósito modelar termodinamicamente o equilíbrio líquido-líquido de sistemas ternários compostos por triglicerídeo, álcool e líquido iônico prótico, avaliando por duas abordagens preditivas de natureza distinta, o UNIFAC-LL e o COSMO-SAC, se esses compostos aumentam ou reduzem a miscibilidade entre o óleo e o álcool. O escopo abrangeu três triglicerídeos com grau crescente de insaturação (trioleína, trilinoleína e trilinolenina), dois álcoois de cadeia curta (etanol e metanol) e dois líquidos iônicos próticos derivados do ânion hexanoato ([2-HEA][Hx] e [DETA][Hx]), o que totaliza doze sistemas ternários avaliados por cada modelo, todos a 298,15~K. A escolha desses dois líquidos iônicos partiu da necessidade de investigar como a estrutura do cátion, hidroxietilamônio monoprotonado em um caso e dietilenotriamônio triprotonado no outro, afeta o comportamento de fases da etapa de separação que sucede o resfriamento do meio reacional da transesterificação alcalina.

A modelagem via UNIFAC-LL exigiu o preenchimento de dezesseis entradas ausentes na parametrização de Magnussen, Rasmussen e Fredenslund (1981), correspondentes a oito pares de grupos considerados nas duas direções: quatro delas envolvem os grupos amino e carboxilato dos líquidos iônicos próticos e entram no cálculo dos doze sistemas centrais, ao passo que as doze restantes envolvem o grupo carboxila e seriam requeridas apenas em sistemas com ácido graxo livre. Adotaram-se, para todas, os valores preditos pelo UNIFAC 2.0 de Hayer, Wendel, Mandt, Hasse e Jirásek (2025), modelo de aprendizado de máquina baseado em completamento de matriz, verificado como compatível com os parâmetros $R_k$ e $Q_k$ de Magnussen, Rasmussen e Fredenslund (1981) para os grupos relevantes a este trabalho. O motor de cálculo, comum aos dois modelos e formado pelo teste de estabilidade por plano tangente, pelo cálculo \textit{flash} por Rachford-Rice e pela localização do ponto crítico, foi validado de forma independente do modelo termodinâmico contra a equação de Margules simétrica, com desvio quadrático médio de $2{,}50 \times 10^{-7}$ em relação à solução analítica. O conjunto formado por motor e modelo foi ainda verificado externamente contra o sistema \{n-hexano + etanol + água\}, com dados de Sørensen e Arlt (1980), obtendo-se RMSD de 0,0872.

Aplicados aos doze sistemas, os dois modelos concordaram quanto aos sistemas com [2-HEA][Hx], nos quais ambos indicaram binodal do tipo 1 para os três triglicerídeos e os dois álcoois, ainda que com discrepância quantitativa apreciável na posição da curva, medida pela distância entre os pontos críticos, entre 0,0616 e 0,3966 em fração molar. Nos sistemas com [DETA][Hx], em contrapartida, os modelos divergiram quanto à própria ocorrência de separação de fases: o COSMO-SAC previu binodal nos seis casos, com área média de 0,4713 do simplex de composições, enquanto no UNIFAC-LL não foi possível construir a curva em nenhum deles. Registra-se esse último resultado nesses termos, e não como predição de miscibilidade completa, uma vez que a ausência de composição instável identificada pelo teste de estabilidade não permite distinguir entre uma fase única estável e uma falha de convergência do procedimento. Os elementos disponíveis sobre a procedência de cada predição foram documentados, a saber, o perfil sigma do [DETA][Hx] verificado quanto à geometria, à estequiometria e à fração polar, de um lado, e os parâmetros de interação obtidos por completamento de matriz aplicado a uma combinação molecular ausente da base de treinamento, de outro; nenhum deles, contudo, constitui verificação do resultado, e os dados deste trabalho não permitem arbitrar a divergência.

O principal achado, e a resposta à pergunta que motivou o estudo, veio da análise da evolução da lacuna de imiscibilidade com a fração de líquido iônico prótico. Nos dezoito sistemas para os quais se obteve curva binodal, sem exceção, a distância entre os dois ramos da curva cresce à medida que o líquido iônico é adicionado, com inclinações entre $+0{,}87$ e $+2{,}95$ em fração molar de triglicerídeo por fração molar de líquido iônico, e a largura da lacuna multiplica-se por fatores de 2,7 a 15,0 entre a menor e a maior faixa amostrada. O mecanismo é comum aos dois modelos: o alargamento decorre quase inteiramente do afastamento do ramo rico em óleo, que avança entre $+0{,}36$ e $+0{,}69$ em fração molar conforme o modelo, enquanto a composição da fase alcoólica permanece praticamente inalterada. Sendo assim, nas condições modeladas, isto é, a 298,15~K e em frações molares de líquido iônico entre 0,047 e 0,621, a hipótese de cossolvência não se confirma: o líquido iônico prótico afasta as fases coexistentes em vez de aproximá-las, comportando-se como agente de separação. Convém destacar que essa conclusão não depende de arbitrar entre os dois modelos, visto que ambos apontam no mesmo sentido, ainda que com intensidades distintas.

A interpretação físico-química proposta atribui esse comportamento à combinação de três mecanismos. A rede de ligações de hidrogênio entre o próton lábil do cátion amônio e a hidroxila do álcool ancora o líquido iônico prótico à fase alcoólica, o que se reflete no coeficiente de distribuição entre as fases, entre 7,3 e 18,4 nos sistemas com [2-HEA][Hx]. Ancorado ali, o líquido iônico compete com o triglicerídeo pela hidroxila disponível e reduz a capacidade solvente dessa fase, competição tanto mais intensa quanto maior a densidade de sítios protonados do cátion, 34\% superior no [DETA][Hx] em relação ao [2-HEA][Hx], o que é coerente com a inclinação média mais acentuada registrada para o primeiro ($+2{,}32$ contra $+1{,}67$ no COSMO-SAC). Por fim, o impedimento estérico das três cadeias esterificadas ao glicerol, quantificado pela diferença de quase três vezes no volume da cavidade COSMO entre a trioleína e uma cadeia graxa isolada, impede que a afinidade dispersiva da cauda hexanoato compense essa perda.

Esse resultado não contradiz os ganhos de rendimento reportados por Gomez et al. (2019) e Trindade et al. (2019), uma vez que três diferenças de escopo separam as duas observações. A concentração empregada nesses trabalhos, de 2 a 3\% em massa, corresponde a frações molares entre aproximadamente 0,005 e 0,028, portanto inteiramente abaixo da menor faixa amostrada neste trabalho. A temperatura difere, sendo de aproximadamente 60\,°C na condição reacional contra 298,15~K nos diagramas aqui calculados. E a grandeza medida é distinta, uma vez que rendimento de reação e equilíbrio de fases podem variar de forma independente. O comportamento identificado, no qual a fase alcoólica não altera sua composição com a adição do líquido iônico, torna a solubilização em fase homogênea uma explicação pouco provável para os ganhos observados e favorece, por eliminação, mecanismos de natureza interfacial.

Do ponto de vista da contribuição, a comparação sistemática entre duas abordagens preditivas de natureza distinta mostra-se particularmente útil em um contexto no qual dados experimentais de equilíbrio líquido-líquido para sistemas com líquidos iônicos próticos e triglicerídeos são inexistentes, de modo que a concordância ou a divergência entre modelos funciona como critério adicional de leitura na falta de validação direta. Os resultados fornecem, ainda, uma base termodinâmica para o dimensionamento da etapa de separação de fases posterior ao resfriamento do meio reacional, ao indicar que a presença residual de líquido iônico prótico nessa etapa tende a favorecer, e não a dificultar, a separação entre as fases.

Como limitações, destacam-se sete pontos discutidos ao longo do Capítulo~\ref{cap:resultados}. Primeiro, a ausência de dados experimentais de equilíbrio líquido-líquido para os sistemas de interesse restringiu toda a avaliação a critérios indiretos, sem verificação contra medida direta. Segundo, a geometria da trilinoleína registrada no banco de dados interno do JCOSMO corresponde a uma conformação enovelada, identificada por inspeção dos arquivos de saída do GAMESS, o que compromete a magnitude dos resultados dessa molécula no COSMO-SAC. Terceiro, os parâmetros de interação do UNIFAC-LL para os grupos dos líquidos iônicos próticos provêm de completamento de matriz aplicado fora do domínio molecular de interesse, e não de calibração direta. Quarto, o par de interação CCOO/COO mostrou-se sensível ao valor adotado, a ponto de a substituição por um valor alternativo eliminar a formação de curva binodal em todos os sistemas, e a determinação de qual valor é fisicamente mais adequado não foi realizada. Quinto, a extrapolação da largura da lacuna para fração nula de líquido iônico fornece valores negativos em doze dos dezoito sistemas, o que implicaria par binário \{triglicerídeo + etanol\} miscível, em desacordo com o comportamento experimental conhecido e coerente com o coeficiente de atividade em diluição infinita de 2,61 obtido para a trioleína em etanol; a aresta binária não pode, por conseguinte, servir de referência quantitativa absoluta. Sexto, a faixa de concentração modelada situa-se uma a duas ordens de grandeza acima da empregada nos estudos experimentais de referência, de modo que a extensão dos resultados ao regime reacional não é assegurada pelos dados obtidos. Sétimo, e possivelmente o mais restritivo, tanto o UNIFAC-LL quanto o COSMO-SAC descrevem a solução como molecularmente dispersa e não representam a formação de agregados, embora os líquidos iônicos próticos estudados apresentem porção apolar majoritária em área de superfície, de 69,8\% e 75,9\%; caso agregados se formem nos sistemas reais, a solubilização de triglicerídeo em seu interior constituiria mecanismo de cossolvência ao qual os modelos empregados são estruturalmente cegos. Soma-se a essas limitações o custo computacional do COSMO-SAC, que demanda cálculos de estrutura eletrônica para cada espécie e torna sua aplicação menos imediata que a de modelos baseados em contribuição de grupos.

Como sugestões para trabalhos futuros, recomenda-se, em primeiro lugar, a obtenção de dados experimentais de equilíbrio líquido-líquido para ao menos um dos doze sistemas estudados, o que permitiria verificar diretamente ambos os modelos e arbitrar a divergência registrada para o [DETA][Hx]. Recomenda-se, ainda, a extensão dos cálculos à faixa de baixa concentração de líquido iônico, entre 0,005 e 0,047 em fração molar, com malha de composições mais fina, capaz de determinar se a curva binodal alcança ou não a aresta binária e de verificar se o comportamento de separação identificado se mantém no regime empregado na reação; a reotimização da geometria da trilinoleína a partir de conformação inicial distinta; a investigação da hipótese de agregação, por simulação de dinâmica molecular ou por determinação experimental da concentração micelar crítica dos dois líquidos iônicos próticos nos meios estudados; a investigação do par de interação CCOO/COO e de sua influência sobre os resultados; e a extensão do estudo à temperatura reacional e a sistemas que incluam água e glicerol, ampliando a aplicabilidade dos resultados à etapa de purificação do biodiesel.